% Options for packages loaded elsewhere
\PassOptionsToPackage{unicode}{hyperref}
\PassOptionsToPackage{hyphens}{url}
%
\documentclass[
]{article}
\usepackage{amsmath,amssymb}
\usepackage{iftex}
\ifPDFTeX
  \usepackage[T1]{fontenc}
  \usepackage[utf8]{inputenc}
  \usepackage{textcomp} % provide euro and other symbols
\else % if luatex or xetex
  \usepackage{unicode-math} % this also loads fontspec
  \defaultfontfeatures{Scale=MatchLowercase}
  \defaultfontfeatures[\rmfamily]{Ligatures=TeX,Scale=1}
\fi
\usepackage{lmodern}
\ifPDFTeX\else
  % xetex/luatex font selection
\fi
% Use upquote if available, for straight quotes in verbatim environments
\IfFileExists{upquote.sty}{\usepackage{upquote}}{}
\IfFileExists{microtype.sty}{% use microtype if available
  \usepackage[]{microtype}
  \UseMicrotypeSet[protrusion]{basicmath} % disable protrusion for tt fonts
}{}
\makeatletter
\@ifundefined{KOMAClassName}{% if non-KOMA class
  \IfFileExists{parskip.sty}{%
    \usepackage{parskip}
  }{% else
    \setlength{\parindent}{0pt}
    \setlength{\parskip}{6pt plus 2pt minus 1pt}}
}{% if KOMA class
  \KOMAoptions{parskip=half}}
\makeatother
\usepackage{xcolor}
\usepackage[margin=1in]{geometry}
\usepackage{color}
\usepackage{fancyvrb}
\newcommand{\VerbBar}{|}
\newcommand{\VERB}{\Verb[commandchars=\\\{\}]}
\DefineVerbatimEnvironment{Highlighting}{Verbatim}{commandchars=\\\{\}}
% Add ',fontsize=\small' for more characters per line
\usepackage{framed}
\definecolor{shadecolor}{RGB}{248,248,248}
\newenvironment{Shaded}{\begin{snugshade}}{\end{snugshade}}
\newcommand{\AlertTok}[1]{\textcolor[rgb]{0.94,0.16,0.16}{#1}}
\newcommand{\AnnotationTok}[1]{\textcolor[rgb]{0.56,0.35,0.01}{\textbf{\textit{#1}}}}
\newcommand{\AttributeTok}[1]{\textcolor[rgb]{0.13,0.29,0.53}{#1}}
\newcommand{\BaseNTok}[1]{\textcolor[rgb]{0.00,0.00,0.81}{#1}}
\newcommand{\BuiltInTok}[1]{#1}
\newcommand{\CharTok}[1]{\textcolor[rgb]{0.31,0.60,0.02}{#1}}
\newcommand{\CommentTok}[1]{\textcolor[rgb]{0.56,0.35,0.01}{\textit{#1}}}
\newcommand{\CommentVarTok}[1]{\textcolor[rgb]{0.56,0.35,0.01}{\textbf{\textit{#1}}}}
\newcommand{\ConstantTok}[1]{\textcolor[rgb]{0.56,0.35,0.01}{#1}}
\newcommand{\ControlFlowTok}[1]{\textcolor[rgb]{0.13,0.29,0.53}{\textbf{#1}}}
\newcommand{\DataTypeTok}[1]{\textcolor[rgb]{0.13,0.29,0.53}{#1}}
\newcommand{\DecValTok}[1]{\textcolor[rgb]{0.00,0.00,0.81}{#1}}
\newcommand{\DocumentationTok}[1]{\textcolor[rgb]{0.56,0.35,0.01}{\textbf{\textit{#1}}}}
\newcommand{\ErrorTok}[1]{\textcolor[rgb]{0.64,0.00,0.00}{\textbf{#1}}}
\newcommand{\ExtensionTok}[1]{#1}
\newcommand{\FloatTok}[1]{\textcolor[rgb]{0.00,0.00,0.81}{#1}}
\newcommand{\FunctionTok}[1]{\textcolor[rgb]{0.13,0.29,0.53}{\textbf{#1}}}
\newcommand{\ImportTok}[1]{#1}
\newcommand{\InformationTok}[1]{\textcolor[rgb]{0.56,0.35,0.01}{\textbf{\textit{#1}}}}
\newcommand{\KeywordTok}[1]{\textcolor[rgb]{0.13,0.29,0.53}{\textbf{#1}}}
\newcommand{\NormalTok}[1]{#1}
\newcommand{\OperatorTok}[1]{\textcolor[rgb]{0.81,0.36,0.00}{\textbf{#1}}}
\newcommand{\OtherTok}[1]{\textcolor[rgb]{0.56,0.35,0.01}{#1}}
\newcommand{\PreprocessorTok}[1]{\textcolor[rgb]{0.56,0.35,0.01}{\textit{#1}}}
\newcommand{\RegionMarkerTok}[1]{#1}
\newcommand{\SpecialCharTok}[1]{\textcolor[rgb]{0.81,0.36,0.00}{\textbf{#1}}}
\newcommand{\SpecialStringTok}[1]{\textcolor[rgb]{0.31,0.60,0.02}{#1}}
\newcommand{\StringTok}[1]{\textcolor[rgb]{0.31,0.60,0.02}{#1}}
\newcommand{\VariableTok}[1]{\textcolor[rgb]{0.00,0.00,0.00}{#1}}
\newcommand{\VerbatimStringTok}[1]{\textcolor[rgb]{0.31,0.60,0.02}{#1}}
\newcommand{\WarningTok}[1]{\textcolor[rgb]{0.56,0.35,0.01}{\textbf{\textit{#1}}}}
\usepackage{longtable,booktabs,array}
\usepackage{calc} % for calculating minipage widths
% Correct order of tables after \paragraph or \subparagraph
\usepackage{etoolbox}
\makeatletter
\patchcmd\longtable{\par}{\if@noskipsec\mbox{}\fi\par}{}{}
\makeatother
% Allow footnotes in longtable head/foot
\IfFileExists{footnotehyper.sty}{\usepackage{footnotehyper}}{\usepackage{footnote}}
\makesavenoteenv{longtable}
\usepackage{graphicx}
\makeatletter
\def\maxwidth{\ifdim\Gin@nat@width>\linewidth\linewidth\else\Gin@nat@width\fi}
\def\maxheight{\ifdim\Gin@nat@height>\textheight\textheight\else\Gin@nat@height\fi}
\makeatother
% Scale images if necessary, so that they will not overflow the page
% margins by default, and it is still possible to overwrite the defaults
% using explicit options in \includegraphics[width, height, ...]{}
\setkeys{Gin}{width=\maxwidth,height=\maxheight,keepaspectratio}
% Set default figure placement to htbp
\makeatletter
\def\fps@figure{htbp}
\makeatother
\setlength{\emergencystretch}{3em} % prevent overfull lines
\providecommand{\tightlist}{%
  \setlength{\itemsep}{0pt}\setlength{\parskip}{0pt}}
\setcounter{secnumdepth}{-\maxdimen} % remove section numbering
\ifLuaTeX
  \usepackage{selnolig}  % disable illegal ligatures
\fi
\IfFileExists{bookmark.sty}{\usepackage{bookmark}}{\usepackage{hyperref}}
\IfFileExists{xurl.sty}{\usepackage{xurl}}{} % add URL line breaks if available
\urlstyle{same}
\hypersetup{
  hidelinks,
  pdfcreator={LaTeX via pandoc}}

\author{}
\date{\vspace{-2.5em}}

\begin{document}

\hypertarget{benchmarking-memo-2-two-measures-of-subnational-democracy-calculation-interpretation-and-use}{%
\section{Benchmarking Memo 2: Two Measures of Subnational Democracy --
Calculation, Interpretation, and
Use}\label{benchmarking-memo-2-two-measures-of-subnational-democracy-calculation-interpretation-and-use}}

\textbf{Date:} July 2, 2026 \textbf{Prepared by:} PM, with assistance
from Claude Code Sonnet 5 \textbf{For:} Michael Coppedge and the
SNVDEM-COL team

\begin{center}\rule{0.5\linewidth}{0.5pt}\end{center}

This memo follows directly from Michael's reply to the June 26
benchmarking memo. That memo (attached) laid out a scale-mismatch
problem: \texttt{EL\_col\_mt} (elections, V-Dem's latent scale) and
\texttt{CL\_col\_mt} (civil liberties, {[}0,1{]} scale) couldn't be
averaged into \texttt{sndem\_final} without distorting the trend.
Michael's reply resolved the open questions -- benchmarked scores are
the intended final measure for all uses, and the two dimensions should
be combined by converting \texttt{CL\_col\_mt} to a Z-score with
\texttt{qnorm()}.

This memo implements that and lays out where things stand: we now have
two complete, validated measures of Colombian subnational democracy --
an \textbf{unbenchmarked} measure and a \textbf{benchmarked} measure --
built for different purposes. Below: how each is calculated, what each
tells us, and how we suggest presenting both in the paper.

\begin{center}\rule{0.5\linewidth}{0.5pt}\end{center}

\hypertarget{the-unbenchmarked-measure-sndem}{%
\subsection{\texorpdfstring{1. The unbenchmarked measure
(\texttt{sndem})}{1. The unbenchmarked measure (sndem)}}\label{the-unbenchmarked-measure-sndem}}

Each municipality's geographic and structural predictors (rurality,
distance from the capital, violence indicators, population density,
etc.) are standardized against the empirical distribution of \emph{other
Colombian municipalities} (an empirical CDF, computed across the full
2000-2023 panel), then combined using V-Dem expert-derived relevance
weights for elections and civil liberties separately:

\begin{Shaded}
\begin{Highlighting}[]
\CommentTok{\# 07\_weighting/01\_weight\_predictors.R}
\NormalTok{snelect  }\OtherTok{=} \FunctionTok{sum}\NormalTok{(predictor\_k }\SpecialCharTok{*}\NormalTok{ weight\_k) }\SpecialCharTok{/} \FunctionTok{sum}\NormalTok{(weight\_k)   }\CommentTok{\# per municipality{-}year}
\NormalTok{sncivlib }\OtherTok{=} \FunctionTok{sum}\NormalTok{(predictor\_k }\SpecialCharTok{*}\NormalTok{ weight\_k) }\SpecialCharTok{/} \FunctionTok{sum}\NormalTok{(weight\_k)   }\CommentTok{\# civil liberties analog}
\NormalTok{sndem    }\OtherTok{=} \FloatTok{0.5} \SpecialCharTok{*}\NormalTok{ (snelect }\SpecialCharTok{+}\NormalTok{ sncivlib)}
\end{Highlighting}
\end{Shaded}

\textbf{What it represents:} each municipality's position \emph{relative
to the rest of Colombia}. It is a within-country, cross-sectional
measure -- well suited to answering ``which kinds of municipalities look
more or less democratic than others in Colombia, and why.'' It reflects
real temporal change to the extent the underlying geopredictors and
expert weights shift year to year, but it is not anchored to any
external, absolute reference point.

\begin{center}\rule{0.5\linewidth}{0.5pt}\end{center}

\hypertarget{the-benchmarked-measure-sndem_final}{%
\subsection{\texorpdfstring{2. The benchmarked measure
(\texttt{sndem\_final})}{2. The benchmarked measure (sndem\_final)}}\label{the-benchmarked-measure-sndem_final}}

\textbf{Step 1 -- map each municipality onto the V-Dem scale.} Each
municipality's unbenchmarked score is re-expressed as a deviation from
the actual V-Dem national estimate for that year, scaled by how much the
country is estimated to vary internally:

\begin{Shaded}
\begin{Highlighting}[]
\CommentTok{\# 08\_benchmark/01\_benchmark.R}
\NormalTok{EL\_col\_mt }\OtherTok{=}\NormalTok{ v2elffelr }\SpecialCharTok{+}\NormalTok{ (snelect  }\SpecialCharTok{{-}}\NormalTok{ snelectyrmean) }\SpecialCharTok{*}\NormalTok{ weighted\_range }\SpecialCharTok{/}\NormalTok{ ELrange\_975\_025}
\NormalTok{CL\_col\_mt }\OtherTok{=}\NormalTok{ CLSNmean  }\SpecialCharTok{+}\NormalTok{ (sncivlib }\SpecialCharTok{{-}}\NormalTok{ CLSNyrmean)   }\SpecialCharTok{*}\NormalTok{ wtdCL\_range   }\SpecialCharTok{/}\NormalTok{ CLrange\_975\_025}
\end{Highlighting}
\end{Shaded}

Precisely, the two municipal-level anchors are built differently, which
matters for how to read them:

\begin{longtable}[]{@{}
  >{\raggedright\arraybackslash}p{(\columnwidth - 4\tabcolsep) * \real{0.3333}}
  >{\raggedright\arraybackslash}p{(\columnwidth - 4\tabcolsep) * \real{0.3333}}
  >{\raggedright\arraybackslash}p{(\columnwidth - 4\tabcolsep) * \real{0.3333}}@{}}
\toprule\noalign{}
\begin{minipage}[b]{\linewidth}\raggedright
Dimension
\end{minipage} & \begin{minipage}[b]{\linewidth}\raggedright
Municipal anchor
\end{minipage} & \begin{minipage}[b]{\linewidth}\raggedright
Built from
\end{minipage} \\
\midrule\noalign{}
\endhead
\bottomrule\noalign{}
\endlastfoot
Elections & \texttt{v2elffelr} & V-Dem's ``Subnational elections free
and fair'' -- latent scale, \textasciitilde{[}-3.5, 3.5{]} \\
Civil liberties &
\texttt{CLSNmean\ =\ v2x\_civlib\ *\ (100\ -\ v2clsnlpct)\ /\ 100} &
V-Dem's Civil Liberties Index (\texttt{v2x\_civlib}, {[}0,1{]}) adjusted
downward by \texttt{v2clsnlpct} -- V-Dem's own estimate of the percent
of the population living under weaker civil liberties than the national
figure \\
\end{longtable}

\texttt{v2elffelr} is used as-is. \texttt{CLSNmean} is not the raw civil
liberties index -- it's \texttt{v2x\_civlib} discounted by how much of
the population V-Dem's own coders judge to be experiencing
weaker-than-national civil liberties.

\texttt{weighted\_range}/\texttt{wtdCL\_range} are
coder-proportion-weighted estimates of how unevenly V-Dem experts judge
Colombia's subnational units to vary:

\begin{Shaded}
\begin{Highlighting}[]
\CommentTok{\# 06\_vdem\_data/01\_weighting.R}
\NormalTok{weighted\_range }\OtherTok{=}\NormalTok{ (snlsff\_1 }\SpecialCharTok{*}\NormalTok{ HPD }\SpecialCharTok{+}\NormalTok{ snlsff\_2 }\SpecialCharTok{*} \DecValTok{2} \SpecialCharTok{*}\NormalTok{ HPD) }\SpecialCharTok{/}\NormalTok{ (snlsff\_0 }\SpecialCharTok{+}\NormalTok{ snlsff\_1 }\SpecialCharTok{+}\NormalTok{ snlsff\_2)}
\end{Highlighting}
\end{Shaded}

(\texttt{snlsff\_0/1/2} = counts of V-Dem coders judging Colombia's
subnational elections ``the same,'' ``somewhat different,'' or
``significantly different''; \texttt{HPD} = the width of V-Dem's
uncertainty interval around \texttt{v2elffelr}. \texttt{wtdCL\_range} is
constructed the same way for civil liberties.)

\textbf{Step 2 -- combine elections and civil liberties on a common,
globally comparable scale.} Per Michael's guidance, \texttt{CL\_col\_mt}
is converted to a Z-score with \texttt{qnorm()}. Both dimensions are
then standardized against fixed constants computed once from the full
V-Dem country-year panel (all countries, 2000-2023) -- not from
Colombia's own data -- so the scale means the same thing regardless of
which country is being scored, which matters once this pipeline is
extended beyond Colombia:

\begin{Shaded}
\begin{Highlighting}[]
\NormalTok{v15\_global      }\OtherTok{\textless{}{-}}\NormalTok{ vdemdata}\SpecialCharTok{::}\NormalTok{vdem }\SpecialCharTok{\%\textgreater{}\%} \FunctionTok{filter}\NormalTok{(year }\SpecialCharTok{\textgreater{}=} \DecValTok{2000}\NormalTok{, year }\SpecialCharTok{\textless{}=} \DecValTok{2023}\NormalTok{)}
\NormalTok{EL\_global\_mean  }\OtherTok{\textless{}{-}} \FunctionTok{mean}\NormalTok{(v15\_global}\SpecialCharTok{$}\NormalTok{v2elffelr, }\AttributeTok{na.rm =} \ConstantTok{TRUE}\NormalTok{)}
\NormalTok{EL\_global\_sd    }\OtherTok{\textless{}{-}} \FunctionTok{sd}\NormalTok{(v15\_global}\SpecialCharTok{$}\NormalTok{v2elffelr,   }\AttributeTok{na.rm =} \ConstantTok{TRUE}\NormalTok{)}

\CommentTok{\# CL global reference must be built the same way as the CL municipal anchor (CLSNmean, i.e.}
\CommentTok{\# v2x\_civlib discounted by v2clsnlpct) or the two aren\textquotesingle{}t comparable {-}{-} see note below.}
\NormalTok{v15\_global      }\OtherTok{\textless{}{-}}\NormalTok{ v15\_global }\SpecialCharTok{\%\textgreater{}\%} \FunctionTok{mutate}\NormalTok{(}\AttributeTok{CLSNmean\_global =}\NormalTok{ v2x\_civlib }\SpecialCharTok{*}\NormalTok{ (}\DecValTok{100} \SpecialCharTok{{-}}\NormalTok{ v2clsnlpct) }\SpecialCharTok{/} \DecValTok{100}\NormalTok{)}
\NormalTok{CLSNmean\_valid  }\OtherTok{\textless{}{-}}\NormalTok{ v15\_global}\SpecialCharTok{$}\NormalTok{CLSNmean\_global[}\SpecialCharTok{!}\FunctionTok{is.na}\NormalTok{(v15\_global}\SpecialCharTok{$}\NormalTok{CLSNmean\_global) }\SpecialCharTok{\&}
\NormalTok{                                                 v15\_global}\SpecialCharTok{$}\NormalTok{CLSNmean\_global }\SpecialCharTok{\textgreater{}} \DecValTok{0} \SpecialCharTok{\&}
\NormalTok{                                                 v15\_global}\SpecialCharTok{$}\NormalTok{CLSNmean\_global }\SpecialCharTok{\textless{}} \DecValTok{1}\NormalTok{]}
\NormalTok{CLz\_global      }\OtherTok{\textless{}{-}} \FunctionTok{qnorm}\NormalTok{(CLSNmean\_valid)}
\NormalTok{CL\_global\_mean  }\OtherTok{\textless{}{-}} \FunctionTok{mean}\NormalTok{(CLz\_global, }\AttributeTok{na.rm =} \ConstantTok{TRUE}\NormalTok{)}
\NormalTok{CL\_global\_sd    }\OtherTok{\textless{}{-}} \FunctionTok{sd}\NormalTok{(CLz\_global,   }\AttributeTok{na.rm =} \ConstantTok{TRUE}\NormalTok{)}

\NormalTok{CL\_col\_z    }\OtherTok{\textless{}{-}} \FunctionTok{qnorm}\NormalTok{(CL\_col\_mt)}
\NormalTok{EL\_col\_gz   }\OtherTok{\textless{}{-}}\NormalTok{ (EL\_col\_mt }\SpecialCharTok{{-}}\NormalTok{ EL\_global\_mean) }\SpecialCharTok{/}\NormalTok{ EL\_global\_sd}
\NormalTok{CL\_col\_gz   }\OtherTok{\textless{}{-}}\NormalTok{ (CL\_col\_z  }\SpecialCharTok{{-}}\NormalTok{ CL\_global\_mean) }\SpecialCharTok{/}\NormalTok{ CL\_global\_sd}
\NormalTok{sndem\_final }\OtherTok{\textless{}{-}} \FloatTok{0.5} \SpecialCharTok{*}\NormalTok{ (EL\_col\_gz }\SpecialCharTok{+}\NormalTok{ CL\_col\_gz)}
\end{Highlighting}
\end{Shaded}

\textbf{A scale asymmetry, found and fixed 2026-07-02.} An earlier
version built \texttt{CL\_global\_mean}/\texttt{CL\_global\_sd} from the
\emph{raw, unadjusted} \texttt{v2x\_civlib}, while Colombia's
municipalities are anchored to \texttt{CLSNmean} -- the same index
discounted by \texttt{v2clsnlpct} (Section 2 above). Elections was
apples-to-apples (\texttt{v2elffelr} used identically at both the
municipal and global level); civil liberties was not. Since
\texttt{v2clsnlpct\ \textgreater{}\ 0} for virtually every country
(global mean \textasciitilde37\%), that mismatch alone pulled Colombia's
\texttt{CL\_col\_gz} down by about \textbf{0.94 global standard
deviations}, independent of any real difference in civil liberties --
enough to flip the substantive story from ``Colombia's civil liberties
lag badly behind the global average'' to ``roughly typical.'' Fixed by
building the global reference the same discounted way as the municipal
anchor (code above). All figures and numbers in this memo reflect the
corrected version.

\textbf{What it represents:} each municipality's position relative to
the \emph{entire V-Dem world}, not just the rest of Colombia. On this
scale, 0 = the average country-year in the full 2000-2023 V-Dem panel,
and units are standard deviations of that global distribution. It is
directly anchored to Colombia's real national trajectory
(\texttt{v2elffelr}, \texttt{v2x\_civlib}) and -- because the
standardization constants are fixed rather than recomputed per country
-- is built to support comparison with other countries once they're
added to the pipeline.

\begin{center}\rule{0.5\linewidth}{0.5pt}\end{center}

\hypertarget{validation}{%
\subsection{3. Validation}\label{validation}}

Three checks confirm the benchmarked measure is doing what it's supposed
to:

\textbf{It tracks Colombia's real national trajectory.} The year-by-year
municipal mean of \texttt{EL\_col\_gz} tracks the (globally
standardized) raw \texttt{v2elffelr} exactly -- by construction, since
\texttt{EL\_col\_mt}'s municipal deviations average to zero within any
year regardless of \texttt{weighted\_range}. \texttt{CL\_col\_gz} tracks
\texttt{qnorm(CLSNmean)} closely, with a small residual gap (expected:
\texttt{mean(qnorm(x))\ !=\ qnorm(mean(x))} for a nonlinear function
applied municipality-by-municipality before averaging, vs.~once to the
national aggregate -- not a data issue).

\begin{figure}
\centering
\includegraphics{../trend_vs_national_anchors.png}
\caption{Municipal means vs.~national anchors}
\end{figure}

\textbf{It shows real, meaningful spatial and temporal variation.} A
single fixed-color-scale map across all 24 years (one ggplot object, so
the scale doesn't reset year to year), using the same palette as the
unbenchmarked facet maps (viridis ``magma''), with limits fixed to the
full panel's own range:

\begin{figure}
\centering
\includegraphics{../sndem_final_facet_allyears.png}
\caption{Subnational democracy, 2000-2023}
\end{figure}

Colombia's national trajectory moves from modestly below the global
average in 2000-2003 (under 0.4 global SD), through roughly the global
average around 2004-2009, to modestly above it from 2010 on (peaking
around 0.35-0.4 SD in 2015-2018), with a dip around 2019-2021 that stays
close to average rather than reversing, and partial recovery by 2023.
Within every year, municipalities visibly differ from one another, not
just over time.

\textbf{Elections and civil liberties, compared directly.} Same palette
applied to both dimensions individually (limits fixed across all three
columns so neither is clipped), so they're comparable to each other and
to the combined measure, at milestone years:

\begin{figure}
\centering
\includegraphics{../sndem_dims_milestone_years.png}
\caption{Electoral vs.~civil liberties vs.~combined}
\end{figure}

Elections moves further above the global average by 2016/2023 than civil
liberties does -- a real, modest difference between the two dimensions,
not an artifact of the color scale.

\textbf{Civil liberties specifically, in context.} Because civil
liberties reads visibly darker than elections in the maps above, it's
worth showing both what's actually happening to Colombia's civil
liberties on their own terms, and how that translates once placed on the
global scale:

\begin{figure}
\centering
\includegraphics{../civil_liberties_in_context.png}
\caption{Civil liberties in context}
\end{figure}

Colombia's own civil liberties level (\texttt{CLSNmean}, top panel)
rises fairly steadily from about 0.38 in 2000 to a peak around 0.59 in
2016-2017, dips after 2018, and partially recovers by 2023 -- a real,
moderate improvement over the period, not a flat or declining line. Once
placed on the global scale (bottom panel), the same trajectory moves
from modestly below the global average (2000-2002) to modestly above it
(2010-2018), matching the elections story in shape, just at a smaller
magnitude. The darker map coloring in early years and the generally
cooler tone of the civil-liberties column relative to elections both
reflect this real, moderate gap -- Colombia's civil liberties have
consistently tracked a bit closer to the world average than its
elections have, not a data or rendering problem.

\begin{center}\rule{0.5\linewidth}{0.5pt}\end{center}

\hypertarget{what-each-measure-is-for}{%
\subsection{4. What each measure is
for}\label{what-each-measure-is-for}}

\textbf{Unbenchmarked (\texttt{sndem}, \texttt{snelect},
\texttt{sncivlib}):} the right choice for within-Colombia analysis --
regressions or maps asking what predicts \emph{relative} democratic
quality across Colombian municipalities, since it isolates spatial
variation from any national-level shift.

\textbf{Benchmarked (\texttt{sndem\_final}, \texttt{EL\_col\_mt},
\texttt{CL\_col\_mt}, and the global-standardized
\texttt{EL\_col\_gz}/\texttt{CL\_col\_gz}):} the right choice whenever
the question involves an absolute or externally anchored reference point
-- Colombia's trajectory against its own historical V-Dem trend, its
position relative to the rest of the world, or (once extended) direct
comparison with another country's benchmarked scores. This is also the
version to use for anything framed around the paper's calibration to
V-Dem, since it's anchored to the same national-level estimates V-Dem
itself reports.

Both are legitimate, validated measures of the same underlying
phenomenon, viewed through different reference frames -- using the wrong
one for a given question is the main risk, not that either is
unreliable.

\begin{center}\rule{0.5\linewidth}{0.5pt}\end{center}

\hypertarget{suggested-presentation-in-the-paper}{%
\subsection{5. Suggested presentation in the
paper}\label{suggested-presentation-in-the-paper}}

\begin{itemize}
\tightlist
\item
  Lead with the \textbf{unbenchmarked} maps/analysis for the paper's
  core contribution -- fine-grained subnational spatial variation within
  Colombia, which V-Dem's national-level estimates cannot show. This is
  the paper's unique empirical contribution.
\item
  Use the \textbf{benchmarked} measure to situate Colombia's overall
  trajectory against the V-Dem global scale, and as the basis for any
  cross-national framing or future extension of the pipeline to other
  countries.
\item
  Include a short methods note distinguishing the two explicitly for
  readers, e.g.: \emph{``We report two related measures. The
  unbenchmarked index (\texttt{sndem}) captures each municipality's
  democratic quality relative to other Colombian municipalities. The
  benchmarked index (\texttt{sndem\_final}) anchors these same municipal
  estimates to V-Dem's national-level estimates for Colombia, placing
  them on a scale comparable across countries. The two are not competing
  estimates -- they answer different questions.''}
\end{itemize}

\begin{center}\rule{0.5\linewidth}{0.5pt}\end{center}

\hypertarget{files}{%
\subsection{Files}\label{files}}

\begin{itemize}
\tightlist
\item
  \texttt{06\_vdem\_data/01\_weighting.R} -- geopredictor weighting and
  national range calculation
\item
  \texttt{07\_weighting/01\_weight\_predictors.R} -- unbenchmarked index
  (\texttt{sndem})
\item
  \texttt{08\_benchmark/01\_benchmark.R} -- benchmarked index
  (\texttt{sndem\_final})
\item
  \texttt{08\_benchmark/02\_validate\_map.R} -- maps in this memo
\item
  \texttt{08\_benchmark/03\_trend\_diagnostics.R} -- national-anchor
  validation
\item
  \texttt{08\_benchmark/memo/benchmarking\_memo.md} -- the June 26 memo
  this one follows up on
\end{itemize}

\end{document}
